\documentclass[12pt,a4paper]{article}

\usepackage{amsmath, amssymb, amsthm}
\usepackage{mathtools}
\usepackage{algorithm}
\usepackage{algpseudocode}
\usepackage{booktabs}
\usepackage{hyperref}
\usepackage{xcolor}
\usepackage{geometry}
\usepackage{graphicx}
\usepackage{enumitem}
\usepackage{microtype}
\usepackage{natbib}
\usepackage{thmtools}

\geometry{margin=1in}

\hypersetup{
    colorlinks=true,
    linkcolor=blue!70!black,
    citecolor=green!50!black,
    urlcolor=cyan!60!black,
}

\theoremstyle{definition}
\newtheorem{definition}{Definition}[section]
\newtheorem{remark}{Remark}[section]

\newcommand{\E}{\mathbb{E}}
\newcommand{\R}{\mathbb{R}}
\newcommand{\piold}{\pi_{\theta_k}}
\newcommand{\pinew}{\pi_{\theta}}
\newcommand{\KL}{D_{\mathrm{KL}}}
\newcommand{\Adv}{A}
\newcommand{\Jcost}{J^C}
\newcommand{\Jrew}{J^R}
\newcommand{\climit}{d}
\newcommand{\lagm}{\lambda}
\newcommand{\Fmat}{\mathbf{F}}
\newcommand{\bvec}{\mathbf{b}}
\newcommand{\gvec}{\mathbf{g}}
\newcommand{\xvec}{\mathbf{x}}
\newcommand{\pvec}{\mathbf{p}}

\title{%
  \textbf{On-Policy Safe Reinforcement Learning Algorithms}\\[0.5em]
  \large A Mathematical Reference for the OmniSafe Implementation
}
\author{%
  OmniSafe Research Reference\\
  \small Based on source code at \texttt{omnisafe/algorithms/on\_policy/}
}
\date{2026}

\begin{document}

\maketitle
\tableofcontents
\newpage

%%%%%%%%%%%%%%%%%%%%%%%%%%%%%%%%%%%%%%%%%%%%%%%%%%%%%%%%%%%%%
\section{Background and Problem Formulation}
\label{sec:background}
%%%%%%%%%%%%%%%%%%%%%%%%%%%%%%%%%%%%%%%%%%%%%%%%%%%%%%%%%%%%%

\subsection{Constrained Markov Decision Process}

Safe reinforcement learning is formulated as a \emph{Constrained Markov Decision Process}
(CMDP)~\citep{altman1999constrained}. A CMDP is a tuple
$(\mathcal{S}, \mathcal{A}, P, r, c, \gamma, \mu_0)$ where:

\begin{itemize}[nosep]
  \item $\mathcal{S}$: state space; $\mathcal{A}$: action space;
  \item $P(s'|s,a)$: transition kernel;
  \item $r(s,a) \in \R$: reward function; $c(s,a) \in \R$: cost function;
  \item $\gamma \in (0,1)$: discount factor; $\mu_0$: initial state distribution.
\end{itemize}

A stochastic policy $\pi : \mathcal{S} \to \Delta(\mathcal{A})$ induces the
\emph{reward return} and \emph{cost return}:
\begin{equation}
  J^R(\pi) = \E_{\tau \sim \pi}\!\left[\sum_{t=0}^{\infty} \gamma^t r(s_t, a_t)\right],
  \qquad
  J^C(\pi) = \E_{\tau \sim \pi}\!\left[\sum_{t=0}^{\infty} \gamma^t c(s_t, a_t)\right].
\end{equation}

The \textbf{safe policy optimization} problem is:
\begin{equation}
  \label{eq:cmdp}
  \max_{\pi} \; J^R(\pi) \qquad \text{subject to} \quad J^C(\pi) \leq d,
\end{equation}
where $d \geq 0$ is a pre-specified cost limit.

\subsection{Value Functions and Advantage Estimation}

For a policy $\pi$ and signal $f \in \{r, c\}$, we define the \emph{value function}
and \emph{action-value function}:
\begin{equation}
  V_f^\pi(s) = \E_{\tau \sim \pi}\!\left[\sum_{t=0}^{\infty}\gamma^t f(s_t,a_t)\;\Big|\;s_0=s\right],
  \quad
  Q_f^\pi(s,a) = f(s,a) + \gamma \E_{s'}[V_f^\pi(s')].
\end{equation}
The \emph{advantage function} is $A_f^\pi(s,a) = Q_f^\pi(s,a) - V_f^\pi(s)$.

\subsubsection{Generalized Advantage Estimation (GAE)}

OmniSafe uses \emph{Generalized Advantage Estimation}~\citep{schulman2015high} to reduce
variance in the policy gradient. Let $\delta_t^f = f(s_t,a_t) + \gamma V_f(s_{t+1}) - V_f(s_t)$
be the temporal-difference residual. The GAE estimate is:
\begin{equation}
  \label{eq:gae}
  \hat{A}_t^{\mathrm{GAE}(\gamma,\lambda)} = \sum_{k=0}^{\infty} (\gamma\lambda)^k \delta_{t+k}^f,
\end{equation}
where $\lambda \in [0,1]$ interpolates between one-step TD ($\lambda=0$, low variance, high
bias) and Monte-Carlo returns ($\lambda=1$, high variance, zero bias).

The corresponding \emph{value target} is:
\begin{equation}
  \hat{V}_f(s_t) = \hat{A}_t^{\mathrm{GAE}} + V_f(s_t).
\end{equation}

In OmniSafe, separate GAE estimates are computed for reward ($\lambda_r$) and cost ($\lambda_c$),
yielding $\hat{A}_t^R$ and $\hat{A}_t^C$ respectively.

\subsubsection{Importance Sampling}

On-policy algorithms update the policy over multiple mini-batches after collecting
data under the \emph{old} policy $\piold$. This requires an importance sampling ratio:
\begin{equation}
  \label{eq:is_ratio}
  \rho_t(\theta) = \frac{\pi_\theta(a_t|s_t)}{\piold(a_t|s_t)} = \exp\!\bigl(\log\pi_\theta(a_t|s_t) - \log\piold(a_t|s_t)\bigr).
\end{equation}

\subsection{Value Function Updates}

All algorithms in OmniSafe update the reward and cost critics by minimizing the MSE loss:
\begin{equation}
  \label{eq:critic_loss}
  \mathcal{L}_V = \frac{1}{N}\sum_{i=1}^N \bigl(\hat{V}_f(s_i) - V_f^\theta(s_i)\bigr)^2 +
  \beta_{\mathrm{norm}}\sum_j \|\theta_j\|^2,
\end{equation}
where the optional $L_2$ regularisation term is controlled by \texttt{use\_critic\_norm} and
\texttt{critic\_norm\_coef} $\beta_{\mathrm{norm}}$.

%%%%%%%%%%%%%%%%%%%%%%%%%%%%%%%%%%%%%%%%%%%%%%%%%%%%%%%%%%%%%
\section{Base Algorithms (No Explicit Safety Constraint)}
\label{sec:base}
%%%%%%%%%%%%%%%%%%%%%%%%%%%%%%%%%%%%%%%%%%%%%%%%%%%%%%%%%%%%%

\subsection{Policy Gradient (PG)}
\label{sec:pg}

\textbf{Reference:}~\citet{sutton1999policy}.

The vanilla policy gradient maximises the expected return by ascending the gradient:
\begin{equation}
  \nabla_\theta J^R = \E_{\tau\sim\piold}\!\left[\sum_t \nabla_\theta\log\pi_\theta(a_t|s_t)\,\hat{A}_t^R\right].
\end{equation}
Using the importance-sampling rewrite, the surrogate loss minimised in OmniSafe is:
\begin{equation}
  \label{eq:pg_loss}
  \mathcal{L}_{\mathrm{PG}}(\theta) = -\E_t\!\left[\rho_t(\theta)\,\hat{A}_t^R\right].
\end{equation}
The advantage surrogate is simply $\hat{A}_t = \hat{A}_t^R$ (cost advantage unused). Gradients
are clipped by $\ell_\infty$-norm \texttt{max\_grad\_norm} before the Adam step. Updates are
repeated for \texttt{update\_iters} passes over mini-batches; early stopping occurs if
$\KL(\piold\|\pi_\theta) > \delta_{\mathrm{KL}}$ (\texttt{target\_kl}).

\subsection{Proximal Policy Optimization (PPO)}
\label{sec:ppo}

\textbf{Reference:}~\citet{schulman2017proximal}.

PPO replaces the un-constrained policy gradient with a clipped surrogate to prevent
destructively large updates:
\begin{equation}
  \label{eq:ppo_loss}
  \mathcal{L}_{\mathrm{PPO}}(\theta) = -\E_t\!\left[
    \min\!\left(\rho_t\hat{A}_t,\;\mathrm{clip}(\rho_t,\,1-\epsilon,\,1+\epsilon)\,\hat{A}_t\right)
  \right]
  - \beta_H \,\mathcal{H}[\pi_\theta],
\end{equation}
where $\epsilon$ (\texttt{clip} $= 0.2$ by default) is the clipping radius and $\beta_H$
(\texttt{entropy\_coef}) weights the entropy bonus. The clipping operation enforces an implicit
trust region: if the ratio deviates too far from 1, its gradient is zeroed.

Compared to PG, PPO inherits the same critic updates (Eq.~\ref{eq:critic_loss}) and the
same advantage surrogate ($\hat{A}_t = \hat{A}_t^R$).

\subsection{Natural Policy Gradient (NPG)}
\label{sec:npg}

\textbf{Reference:}~\citet{kakade2001natural}.

The natural gradient pre-conditions the ordinary gradient by the inverse Fisher information
matrix $\Fmat_\theta^{-1}$:
\begin{equation}
  \tilde{\nabla}_\theta J^R = \Fmat_\theta^{-1}\,\nabla_\theta J^R,
  \qquad
  \Fmat_\theta = \E_{\piold}\!\left[\nabla_\theta\log\pi_\theta\,(\nabla_\theta\log\pi_\theta)^\top\right].
\end{equation}
Computing $\Fmat^{-1}$ directly is infeasible for neural networks.  Instead, OmniSafe
approximates the \emph{Fisher-vector product} (FVP) $\Fmat_\theta\,\mathbf{v}$ via
automatic differentiation of the KL divergence:
\begin{equation}
  \label{eq:fvp}
  \Fmat_\theta\,\mathbf{v}
  = \nabla_\theta\!\bigl[(\nabla_\theta \KL(\piold \| \pi_\theta))^\top\mathbf{v}\bigr]
  + \alpha_{\mathrm{damp}}\,\mathbf{v},
\end{equation}
where $\alpha_{\mathrm{damp}}$ (\texttt{cg\_damping}) is a Tikhonov regulariser for numerical
stability.

The natural gradient direction $\xvec = \Fmat_\theta^{-1}\gvec$ (where
$\gvec = -\nabla_\theta \mathcal{L}_{\mathrm{PG}}$) is obtained by solving the linear system
$\Fmat\,\xvec = \gvec$ with the \emph{conjugate gradient} (CG) algorithm
(\texttt{cg\_iters} $= 15$ iterations).

The step size $\alpha$ is set so that the KL divergence equals the trust-region budget:
\begin{equation}
  \label{eq:npg_step}
  \alpha = \sqrt{\frac{2\delta_{\mathrm{KL}}}{\xvec^\top\Fmat\xvec}},
  \qquad
  \theta_{\mathrm{new}} = \theta_{\mathrm{old}} + \alpha\,\xvec.
\end{equation}
NPG does \emph{not} perform line search; it takes the full step unconditionally.

\begin{algorithm}[h]
\caption{Natural Policy Gradient (NPG) -- actor update}
\label{alg:npg}
\begin{algorithmic}[1]
\Require observations $\{s_i\}$, actions $\{a_i\}$, log-probs $\{\log\piold(a_i|s_i)\}$, advantages $\{\hat{A}_i^R\}$
\State Compute policy-gradient: $\gvec \leftarrow -\nabla_\theta\mathcal{L}_{\mathrm{PG}}$
\State Solve $\Fmat\,\xvec = \gvec$ via CG using the FVP oracle (Eq.~\ref{eq:fvp})
\State Compute $\alpha \leftarrow \sqrt{2\delta_{\mathrm{KL}}/(\xvec^\top\Fmat\xvec)}$
\State $\theta \leftarrow \theta + \alpha\,\xvec$
\end{algorithmic}
\end{algorithm}

\subsection{Trust Region Policy Optimization (TRPO)}
\label{sec:trpo}

\textbf{Reference:}~\citet{schulman2015trust}.

TRPO adds a backtracking line search to NPG to ensure both (i) actual improvement in the
surrogate objective and (ii) satisfaction of the KL trust region:
\begin{equation}
  \label{eq:trpo_obj}
  \max_\theta \;\E_t\!\left[\rho_t\hat{A}_t^R\right]
  \quad\text{s.t.}\quad \E_s\!\left[\KL(\piold(\cdot|s)\|\pi_\theta(\cdot|s))\right] \leq \delta_{\mathrm{KL}}.
\end{equation}
Starting from the NPG candidate $\theta' = \theta + \xvec$ (with the full NPG step), TRPO
performs exponential back-off: $\theta'_j = \theta + \beta^j\,\alpha\xvec$ for
$j=0,1,\ldots,J$ (\texttt{update\_iters} $ = 10$, decay $\beta = 0.8$), accepting the
first $j$ for which:
\begin{enumerate}[nosep]
  \item $\Delta J^R = \mathcal{L}_{\mathrm{before}} - \mathcal{L}(\theta'_j) > 0$ (surrogate improves), and
  \item $\KL(\piold\|\pi_{\theta'_j}) \leq \delta_{\mathrm{KL}}$.
\end{enumerate}
If no step satisfies both conditions, the policy is not updated.

%%%%%%%%%%%%%%%%%%%%%%%%%%%%%%%%%%%%%%%%%%%%%%%%%%%%%%%%%%%%%
\section{Na\"{i}ve Lagrangian Methods}
\label{sec:lagrange}
%%%%%%%%%%%%%%%%%%%%%%%%%%%%%%%%%%%%%%%%%%%%%%%%%%%%%%%%%%%%%

\subsection{The Lagrangian Relaxation}

The CMDP (Eq.~\ref{eq:cmdp}) can be converted to an unconstrained saddle-point problem by
introducing a Lagrange multiplier $\lagm \geq 0$:
\begin{equation}
  \max_\pi \min_{\lagm \geq 0}\;J^R(\pi) - \lagm\bigl(J^C(\pi) - d\bigr).
\end{equation}
This gives the \emph{penalised advantage surrogate} used by all na\"{i}ve Lagrangian algorithms:
\begin{equation}
  \label{eq:lag_adv}
  \hat{A}_t^{\mathrm{Lag}} = \frac{\hat{A}_t^R - \lagm\,\hat{A}_t^C}{1 + \lagm}.
\end{equation}
The normalisation by $(1+\lagm)$ is an implementation detail that stabilises training.

\subsection{Lagrange Multiplier Update}

The multiplier $\lagm$ is a differentiable parameter updated by gradient descent on the
dual objective:
\begin{equation}
  \label{eq:lag_update}
  \mathcal{L}_\lagm = -\lagm\,(J^C_{\mathrm{avg}} - d),
  \qquad
  \lagm \leftarrow \mathrm{clip}\!\left(\lagm - \eta_\lagm\,\frac{\partial\mathcal{L}_\lagm}{\partial\lagm},\; 0,\; \lagm_{\max}\right),
\end{equation}
which amounts to $\lagm \leftarrow \mathrm{clip}(\lagm + \eta_\lagm(J^C_{\mathrm{avg}}-d), 0, \lagm_{\max})$.
Here $\eta_\lagm$ (\texttt{lambda\_lr}) is the multiplier learning rate and $\lagm_{\max}$
(\texttt{lagrangian\_upper\_bound}) is an optional upper bound.

\subsection{PPO-Lag}
\label{sec:ppolag}

PPO-Lag~\citep{ray2019benchmarking} substitutes the penalised advantage (Eq.~\ref{eq:lag_adv})
into the PPO clipped objective (Eq.~\ref{eq:ppo_loss}):
\begin{equation}
  \mathcal{L}_{\mathrm{PPO\text{-}Lag}}(\theta) = -\E_t\!\left[
    \min\!\Bigl(\rho_t\hat{A}_t^{\mathrm{Lag}},\;\mathrm{clip}(\rho_t,1\!-\!\epsilon,1\!+\!\epsilon)\hat{A}_t^{\mathrm{Lag}}\Bigr)
  \right].
\end{equation}
Before each policy update, $\lagm$ is updated via Eq.~\ref{eq:lag_update} using the
episodic cost average $J^C_{\mathrm{avg}}$ from the logger.

\subsection{TRPO-Lag}
\label{sec:trpolag}

TRPO-Lag uses the same penalised advantage $\hat{A}_t^{\mathrm{Lag}}$ as the objective but
employs the TRPO natural-gradient update with line search (Section~\ref{sec:trpo}) instead
of the PPO clipped objective. The Lagrange update rule is identical to Eq.~\ref{eq:lag_update}.

\subsection{Reward Constrained Policy Optimization (RCPO)}
\label{sec:rcpo}

\textbf{Reference:}~\citet{tessler2018reward}.

RCPO is conceptually identical to TRPO-Lag but builds on the NPG base class, combining
the natural gradient direction (no line search) with the penalised Lagrangian advantage:
\begin{equation}
  \hat{A}_t^{\mathrm{RCPO}} = \frac{\hat{A}_t^R - \lagm\,\hat{A}_t^C}{1+\lagm}.
\end{equation}
The multiplier update follows Eq.~\ref{eq:lag_update}.

\subsection{Primal-Dual Optimisation (PDO)}
\label{sec:pdo}

PDO~\citep{chow2017risk} is the simplest Lagrangian baseline: a vanilla policy gradient
with the penalised advantage $\hat{A}_t^{\mathrm{Lag}}$ and the gradient-based multiplier
update (Eq.~\ref{eq:lag_update}). No clipping, natural gradient, or line search is applied.

%%%%%%%%%%%%%%%%%%%%%%%%%%%%%%%%%%%%%%%%%%%%%%%%%%%%%%%%%%%%%
\section{Second-Order Constrained Methods}
\label{sec:secondorder}
%%%%%%%%%%%%%%%%%%%%%%%%%%%%%%%%%%%%%%%%%%%%%%%%%%%%%%%%%%%%%

\subsection{Constrained Policy Optimization (CPO)}
\label{sec:cpo}

\textbf{Reference:}~\citet{achiam2017constrained}.

CPO solves, at each iteration, the following constrained surrogate problem:
\begin{equation}
  \label{eq:cpo_problem}
  \max_\theta \;\E_t\!\left[\rho_t\hat{A}_t^R\right]
  \quad\text{s.t.}\quad
  \E_t\!\left[\rho_t\hat{A}_t^C\right] \leq c,
  \quad
  \E_s\!\left[\KL(\piold\|\pi_\theta)\right] \leq \delta_{\mathrm{KL}},
\end{equation}
where $c = J^C_{\mathrm{avg}} - d$ is the current constraint violation (signed).

\subsubsection{Quadratic Approximation}

Using the Fisher matrix $\Fmat$ as a local metric and defining:
\begin{align}
  \gvec &= \nabla_\theta \mathcal{L}_{\mathrm{PG}} \quad \text{(reward gradient)}, \\
  \bvec &= \nabla_\theta \mathcal{L}_{\mathrm{cost}} \quad \text{(cost gradient)},
\end{align}
the problem (Eq.~\ref{eq:cpo_problem}) is approximated as:
\begin{equation}
  \label{eq:cpo_qp}
  \max_{\Delta\theta}\;\gvec^\top\Delta\theta
  \quad\text{s.t.}\quad
  \bvec^\top\Delta\theta + c \leq 0,
  \quad
  \Delta\theta^\top\Fmat\Delta\theta \leq 2\delta_{\mathrm{KL}}.
\end{equation}
Define:
\begin{equation}
  \xvec = \Fmat^{-1}\gvec,\quad \pvec = \Fmat^{-1}\bvec \quad \text{(via CG)},
\end{equation}
and the scalar quantities:
\begin{equation}
  q = \xvec^\top\Fmat\xvec,\quad
  r = \gvec^\top\pvec,\quad
  s = \bvec^\top\pvec,\quad
  A = q - \frac{r^2}{s},\quad
  B = 2\delta_{\mathrm{KL}} - \frac{c^2}{s}.
\end{equation}

\subsubsection{Analytical Solution by Case}

The optimal dual variables $(\lambda^\star, \nu^\star)$ and the step direction depend on
the geometry of the feasible region. OmniSafe implements four cases:

\begin{description}
  \item[Case 4] ($\bvec^\top\bvec \leq \epsilon_b$ and $c < 0$):
    Cost gradient is negligible and current point is feasible. Use the plain TRPO step:
    \begin{equation}
      \Delta\theta = \alpha\,\xvec, \quad \alpha = \sqrt{2\delta_{\mathrm{KL}}/q}, \quad \nu^\star=0,\; \lambda^\star = 1/\alpha.
    \end{equation}

  \item[Case 3] ($c < 0$, $B < 0$):
    The current point is feasible and the cost half-space does not intersect the trust region.
    Same as Case 4.

  \item[Cases 1 \& 2] ($c \geq 0$ or $B \geq 0$):
    The QP has an active cost constraint. The dual objective is:
    \begin{equation}
      f_a(\lambda) = -\tfrac{1}{2}\!\left(\frac{A}{\lambda} + B\lambda\right) - \frac{rc}{s},
      \qquad
      f_b(\lambda) = -\tfrac{1}{2}\!\left(\frac{q}{\lambda} + 2\delta_{\mathrm{KL}}\lambda\right).
    \end{equation}
    The optimal $\lambda^\star$ is:
    \begin{equation}
      \lambda^\star = \argmax\!\left(f_a(\lambda_a^\star),\; f_b(\lambda_b^\star)\right),
    \end{equation}
    where $\lambda_a^\star = \mathrm{clip}(\sqrt{A/B}, \cdot)$ and $\lambda_b^\star = \mathrm{clip}(\sqrt{q/(2\delta_{\mathrm{KL}})}, \cdot)$ are projected into their respective feasible intervals. Then:
    \begin{equation}
      \nu^\star = \frac{\max(0,\;\lambda^\star c - r)}{s},
      \qquad
      \Delta\theta = \frac{1}{\lambda^\star}(\xvec - \nu^\star\pvec).
    \end{equation}

  \item[Case 0] (full trust region infeasible):
    Fall back to purely decreasing cost:
    \begin{equation}
      \Delta\theta = -\nu^\star\pvec, \quad \nu^\star = \sqrt{2\delta_{\mathrm{KL}}/s}.
    \end{equation}
\end{description}

\subsubsection{CPO Line Search}

After computing the analytical direction, CPO applies a backtracking line search with
decay $\beta = 0.8$ (up to 20 steps) checking:
\begin{enumerate}[nosep]
  \item For Cases~3/4: surrogate reward improves.
  \item $J^C(\theta') - J^C(\theta) \leq \max(0, c)$ (cost does not worsen beyond current violation).
  \item $\KL(\piold\|\pi_{\theta'}) \leq \delta_{\mathrm{KL}}$.
\end{enumerate}

\subsection{Projection-Based Constrained Policy Optimization (PCPO)}
\label{sec:pcpo}

\textbf{Reference:}~\citet{yang2020projection}.

PCPO simplifies CPO by replacing the case analysis with a direct two-step projection.
The unconstrained natural gradient step is first computed, then projected onto the
intersection of the KL trust region and the cost half-space:
\begin{equation}
  \label{eq:pcpo_step}
  \Delta\theta = \sqrt{\frac{2\delta_{\mathrm{KL}}}{q}}\,\Fmat^{-1}\gvec
  - \max\!\!\left(0,\;\frac{\sqrt{2\delta_{\mathrm{KL}}/q}\cdot r + c}{s}\right)\Fmat^{-1}\bvec.
\end{equation}
This can be interpreted as: (i) scale the reward gradient to saturate the KL budget, then
(ii) subtract a correction along the cost gradient to satisfy the linear constraint. PCPO
inherits the CPO line search for final step acceptance.

%%%%%%%%%%%%%%%%%%%%%%%%%%%%%%%%%%%%%%%%%%%%%%%%%%%%%%%%%%%%%
\section{First-Order Constrained Methods}
\label{sec:firstorder}
%%%%%%%%%%%%%%%%%%%%%%%%%%%%%%%%%%%%%%%%%%%%%%%%%%%%%%%%%%%%%

\subsection{First Order Constrained Optimisation in Policy Space (FOCOPS)}
\label{sec:focops}

\textbf{Reference:}~\citet{zhang2020first}.

FOCOPS maintains a Lagrange multiplier $\lagm$ and solves, in the space of non-parametric
policies, a KL-penalised objective to find a locally-optimal constrained policy $\pi^*$,
then projects it back to the parametric family.

\subsubsection{Non-Parametric Optimal Policy}

The closed-form solution of the non-parametric inner problem is:
\begin{equation}
  \pi^*(a|s) \propto \piold(a|s)\,\exp\!\left(\frac{1}{\eta}\bigl(A^R_{\piold}(s,a) - \lagm A^C_{\piold}(s,a)\bigr)\right),
\end{equation}
where $\eta > 0$ (\texttt{focops\_lam}) is a temperature parameter.

\subsubsection{Parametric Projection Loss}

The policy network is updated by minimising the KL distance to $\pi^*$, restricted to
data points where the current KL from the old policy is within a threshold $\eta'$
(\texttt{focops\_eta}):
\begin{equation}
  \label{eq:focops_loss}
  \mathcal{L}_{\mathrm{FOCOPS}}(\theta) = \E_t\!\left[
    \left(\KL(\pi_\theta(\cdot|s_t)\|\piold(\cdot|s_t)) - \frac{1}{\eta}\,\rho_t\,\hat{A}_t^{\mathrm{Lag}}\right)
    \cdot\mathbf{1}\!\left[\KL(\pi_\theta\|\piold) \leq \eta'\right]
  \right]
  - \beta_H\mathcal{H}[\pi_\theta],
\end{equation}
where $\hat{A}_t^{\mathrm{Lag}}$ is the penalised advantage from Eq.~\ref{eq:lag_adv}.
The indicator masks updates for data points where the policy has already moved too far
from $\piold$, acting as a soft trust region.

The Lagrange multiplier is updated using Eq.~\ref{eq:lag_update} before the policy update.

\subsection{Constrained Update Projection (CUP)}
\label{sec:cup}

\textbf{Reference:}~\citet{yang2022constrained}.

CUP is a two-stage algorithm. The first stage maximises the reward via a standard PPO
update (Section~\ref{sec:ppo}). The second stage then projects the policy back towards
constraint satisfaction via gradient descent on a cost-penalised KL loss.

\subsubsection{Stage 1: Reward Update}

Identical to PPO (Eq.~\ref{eq:ppo_loss}) with the reward advantage $\hat{A}_t^R$.

\subsubsection{Stage 2: Cost Projection}

The second-stage loss penalises cost advantage while regularising by KL divergence from
the pre-update policy $\pi_{\theta_0}$ (before Stage 1):
\begin{equation}
  \label{eq:cup_cost_loss}
  \mathcal{L}_{\mathrm{CUP\text{-}cost}}(\theta) = \E_t\!\left[
    \lagm\,\frac{1-\gamma\nu}{1-\gamma}\,\rho_t\,\hat{A}_t^C
    + \KL(\pi_\theta(\cdot|s_t)\|\pi_{\theta_0}(\cdot|s_t))
  \right],
\end{equation}
where $\nu = \lambda_c$ is the cost GAE parameter and $\frac{1-\gamma\nu}{1-\gamma}$ is a
normalisation constant derived from the discounted cost return identity. The KL term
prevents the cost projection from moving too far from the reward-updated policy.

Updates iterate for \texttt{update\_iters} steps with early KL stopping.

%%%%%%%%%%%%%%%%%%%%%%%%%%%%%%%%%%%%%%%%%%%%%%%%%%%%%%%%%%%%%
\section{PID Lagrangian Methods}
\label{sec:pid}
%%%%%%%%%%%%%%%%%%%%%%%%%%%%%%%%%%%%%%%%%%%%%%%%%%%%%%%%%%%%%

\subsection{PID Controller for the Lagrange Multiplier}

\textbf{Reference:}~\citet{stooke2020responsive}.

The na\"{i}ve gradient update of $\lagm$ (Eq.~\ref{eq:lag_update}) can be slow to respond
to sudden constraint violations. PID-Lagrangian replaces it with a
Proportional-Integral-Derivative controller that uses the error signal
$e_t = J^C_{\mathrm{avg}} - d$:
\begin{align}
  \text{P term:}\quad &\tilde{e}_t = \alpha_P\,\tilde{e}_{t-1} + (1-\alpha_P)\,e_t, \label{eq:pid_p}\\
  \text{I term:}\quad &I_t = \max(0,\; I_{t-1} + K_i\,e_t), \label{eq:pid_i}\\
  \text{D term:}\quad &\tilde{C}_t = \alpha_D\,\tilde{C}_{t-1} + (1-\alpha_D)\,J^C_{\mathrm{avg}},
  \quad D_t = \max(0,\;\tilde{C}_t - \tilde{C}_{t-\tau}), \label{eq:pid_d}\\
  \text{Output:}\quad &\lagm_t = \max\!\left(0,\; K_p\,\tilde{e}_t + I_t + K_d\,D_t\right), \label{eq:pid_out}
\end{align}
where $\alpha_P$, $\alpha_D$ are EMA smoothing coefficients (\texttt{pid\_delta\_p\_ema\_alpha},
\texttt{pid\_delta\_d\_ema\_alpha}), $K_p$, $K_i$, $K_d$ are the PID gains, and $\tau$
(\texttt{pid\_d\_delay}) is the delay for the derivative window. The integral is
initialised to \texttt{lagrangian\_multiplier\_init} and is clamped to zero from below.

\subsection{TRPO-PID}
\label{sec:trpopid}

TRPO-PID combines TRPO (Section~\ref{sec:trpo}) with the PID multiplier
(Eqs.~\ref{eq:pid_p}--\ref{eq:pid_out}). The penalised advantage is:
\begin{equation}
  \hat{A}_t^{\mathrm{TRPO\text{-}PID}} = \frac{\hat{A}_t^R - \lagm_t\,\hat{A}_t^C}{1+\lagm_t}.
\end{equation}

\subsection{CPPOPID}
\label{sec:cppopid}

CPPOPID replaces the TRPO update with a PPO clipped objective while retaining the PID
multiplier:
\begin{equation}
  \mathcal{L}_{\mathrm{CPPO\text{-}PID}}(\theta) = -\E_t\!\left[
    \min\!\Bigl(\rho_t\hat{A}_t^{\mathrm{PID}},\;\mathrm{clip}(\rho_t,1\!-\!\epsilon,1\!+\!\epsilon)\hat{A}_t^{\mathrm{PID}}\Bigr)
  \right],
  \quad \hat{A}_t^{\mathrm{PID}} = \frac{\hat{A}_t^R - \lagm_t\hat{A}_t^C}{1+\lagm_t}.
\end{equation}

%%%%%%%%%%%%%%%%%%%%%%%%%%%%%%%%%%%%%%%%%%%%%%%%%%%%%%%%%%%%%
\section{Penalty Function Methods}
\label{sec:penalty}
%%%%%%%%%%%%%%%%%%%%%%%%%%%%%%%%%%%%%%%%%%%%%%%%%%%%%%%%%%%%%

\subsection{Penalised Proximal Policy Optimisation (P3O)}
\label{sec:p3o}

\textbf{Reference:}~\citet{liu2022penalized}.

P3O adds a direct ReLU penalty on the cost surrogate to the PPO objective:
\begin{equation}
  \mathcal{L}_{\mathrm{P3O}}(\theta) = \mathcal{L}_{\mathrm{PPO}}(\theta)
  + \kappa\,\mathrm{ReLU}\!\left(\E_t\!\left[\rho_t\hat{A}_t^C\right] + c\right),
\end{equation}
where $c = J^C_{\mathrm{avg}} - d$ is the current signed violation and $\kappa$
(\texttt{kappa}) is a fixed penalty coefficient. Unlike Lagrangian methods, the penalty
coefficient $\kappa$ is not adapted; it is a fixed hyperparameter. The objective is a
single loss combining reward and cost in one backward pass.

\subsection{Interior-Point Policy Optimisation (IPO)}
\label{sec:ipo}

\textbf{Reference:}~\citet{liu2020ipo}.

IPO uses an interior-point (log-barrier) inspired dynamic penalty that diverges as the cost
approaches the limit:
\begin{equation}
  \label{eq:ipo_penalty}
  \kappa_t = \frac{\kappa_0}{d - J^C_{\mathrm{avg}} + \epsilon},
  \quad \kappa_t \leftarrow \min(\kappa_t, \kappa_{\max}),
\end{equation}
where $\kappa_0$ (\texttt{kappa}) scales the penalty, $\epsilon$ is a small constant for
numerical stability, and $\kappa_{\max}$ (\texttt{penalty\_max}) caps the penalty to prevent
training instability. The penalised advantage is:
\begin{equation}
  \hat{A}_t^{\mathrm{IPO}} = \frac{\hat{A}_t^R - \kappa_t\,\hat{A}_t^C}{1+\kappa_t},
\end{equation}
inserted into the standard PPO objective. The penalty automatically increases as $J^C$
approaches $d$, providing interior-point-like barrier behaviour without requiring dual
variable optimisation.

%%%%%%%%%%%%%%%%%%%%%%%%%%%%%%%%%%%%%%%%%%%%%%%%%%%%%%%%%%%%%
\section{Primal Feasibility Method}
\label{sec:primal}
%%%%%%%%%%%%%%%%%%%%%%%%%%%%%%%%%%%%%%%%%%%%%%%%%%%%%%%%%%%%%

\subsection{Constrained Reward Policy Optimisation (OnCRPO)}
\label{sec:crpo}

\textbf{Reference:}~\citet{xu2021crpo}.

CRPO avoids the Lagrangian entirely and switches between two objectives based on the
current constraint satisfaction:
\begin{equation}
  \label{eq:crpo_adv}
  \hat{A}_t^{\mathrm{CRPO}} =
  \begin{cases}
    \hat{A}_t^R & \text{if } J^C_{\mathrm{avg}} \leq d + \Delta, \\
    -\hat{A}_t^C & \text{otherwise,}
  \end{cases}
\end{equation}
where $\Delta$ (\texttt{distance}) is a tolerance slack. When the constraint is satisfied,
the algorithm performs a standard reward-maximisation step; when violated, it performs a
cost-reduction step. The TRPO update (natural gradient + line search) is applied in both
cases. This primal approach provably converges under mild assumptions~\citep{xu2021crpo}.

%%%%%%%%%%%%%%%%%%%%%%%%%%%%%%%%%%%%%%%%%%%%%%%%%%%%%%%%%%%%%
\section{State Augmentation Methods}
\label{sec:stateaug}
%%%%%%%%%%%%%%%%%%%%%%%%%%%%%%%%%%%%%%%%%%%%%%%%%%%%%%%%%%%%%

State augmentation methods modify the \emph{environment} rather than the policy update,
encoding safety information directly into the observation. Standard PG algorithms (PPO,
TRPO) are then applied without modification.

\subsection{Saute RL (PPO-Saute, TRPO-Saute)}
\label{sec:saute}

\textbf{Reference:}~\citet{sootla2022saute}.

Saute augments the state with a running \emph{safety budget} $z_t$:
\begin{equation}
  \tilde{s}_t = (s_t, z_t), \qquad z_{t+1} = \frac{z_t - c(s_t,a_t)}{\gamma},
\end{equation}
with $z_0 = d/(1-\gamma)$ (the per-step budget). The reward is replaced by a safe reward:
\begin{equation}
  \tilde{r}(s_t,a_t,z_t) =
  \begin{cases}
    r(s_t,a_t) & \text{if } z_t > 0, \\
    r_{\mathrm{unsafe}} & \text{otherwise.}
  \end{cases}
\end{equation}
Since the constraint information is embedded in the state, standard unconstrained PPO/TRPO
can be applied to the augmented MDP $(\mathcal{S}\times\R, \mathcal{A}, \tilde{P}, \tilde{r}, \gamma, \tilde{\mu}_0)$.

\subsection{Simmer RL (PPO-Simmer-PID, TRPO-Simmer-PID)}
\label{sec:simmer}

\textbf{Reference:}~\citet{sootla2022simmer}.

Simmer extends Saute by actively controlling the safety budget $z_t$ via a PID controller.
At each epoch, the observed cost $J^C_{\mathrm{avg}}$ is fed into a budget controller that
adjusts the initial budget $z_0$ for the next epoch:
\begin{equation}
  z_0^{(k+1)} = z_0^{(k)} + \mathrm{PID}(J^C_{\mathrm{avg}}^{(k)} - d).
\end{equation}
The PID controller has the same structure as in Eqs.~\ref{eq:pid_p}--\ref{eq:pid_out},
but it modulates the safety budget rather than a penalty multiplier. Within each episode,
the Saute state augmentation is used unchanged.

%%%%%%%%%%%%%%%%%%%%%%%%%%%%%%%%%%%%%%%%%%%%%%%%%%%%%%%%%%%%%
\section{Early Termination}
\label{sec:et}
%%%%%%%%%%%%%%%%%%%%%%%%%%%%%%%%%%%%%%%%%%%%%%%%%%%%%%%%%%%%%

\subsection{PPO-EarlyTerminated and TRPO-EarlyTerminated}
\label{sec:early}

\textbf{Reference:}~\citet{ji2023omnisafe}.

Early termination is the simplest safety mechanism: episodes are forcibly terminated when
the cumulative cost within an episode exceeds a per-episode threshold $d_{\mathrm{ep}}$.
Formally, the \emph{effective horizon} is:
\begin{equation}
  T^\star(\tau) = \min\!\left(T,\; \inf\!\left\{t : \sum_{k=0}^t c(s_k,a_k) > d_{\mathrm{ep}}\right\}\right).
\end{equation}
All returns are computed only up to $T^\star$, so the policy naturally learns to avoid
constraint violations that cut trajectories short. PPO or TRPO is then applied to the
truncated trajectories without any additional constraint-handling modification.

%%%%%%%%%%%%%%%%%%%%%%%%%%%%%%%%%%%%%%%%%%%%%%%%%%%%%%%%%%%%%
\section{Summary and Comparison}
\label{sec:summary}
%%%%%%%%%%%%%%%%%%%%%%%%%%%%%%%%%%%%%%%%%%%%%%%%%%%%%%%%%%%%%

\begin{table}[h]
\centering
\caption{Summary of on-policy safe RL algorithms in OmniSafe.}
\label{tab:summary}
\small
\setlength{\tabcolsep}{4pt}
\begin{tabular}{lllll}
\toprule
\textbf{Algorithm} & \textbf{Category} & \textbf{Base Update} & \textbf{Constraint Handling} & \textbf{Multiplier} \\
\midrule
PG & Unconstrained & 1st-order PG & --- & --- \\
PPO & Unconstrained & Clipped PG & --- & --- \\
NPG & Unconstrained & Natural PG & --- & --- \\
TRPO & Unconstrained & Natural PG + LS & --- & --- \\
\midrule
PDO & Na\"{i}ve Lagrange & 1st-order PG & Penalised adv. & Gradient \\
PPO-Lag & Na\"{i}ve Lagrange & Clipped PG & Penalised adv. & Gradient \\
TRPO-Lag & Na\"{i}ve Lagrange & Natural PG + LS & Penalised adv. & Gradient \\
RCPO & Na\"{i}ve Lagrange & Natural PG & Penalised adv. & Gradient \\
\midrule
CPO & 2nd-order & Natural PG + LS & Dual QP (4 cases) & Analytical \\
PCPO & 2nd-order & Natural PG + LS & Projection & Analytical \\
\midrule
FOCOPS & 1st-order & Gradient + KL clip & KL-reg.\ penalised adv. & Gradient \\
CUP & 1st-order & Clipped PG + cost step & Two-stage KL projection & Gradient \\
\midrule
TRPO-PID & PID Lagrange & Natural PG + LS & Penalised adv. & PID \\
CPPOPID & PID Lagrange & Clipped PG & Penalised adv. & PID \\
\midrule
P3O & Penalty & Clipped PG & ReLU penalty & Fixed \\
IPO & Penalty & Clipped PG & Log-barrier penalty & Dynamic \\
\midrule
OnCRPO & Primal & Natural PG + LS & Objective switching & None \\
\midrule
PPO-Saute & State Aug. & Clipped PG & Augmented state & None \\
TRPO-Saute & State Aug. & Natural PG + LS & Augmented state & None \\
PPO-Simmer-PID & State Aug. & Clipped PG & Aug.\ state + PID budget & PID \\
TRPO-Simmer-PID & State Aug. & Natural PG + LS & Aug.\ state + PID budget & PID \\
\midrule
PPO-EarlyTerm & Early Term. & Clipped PG & Episode termination & None \\
TRPO-EarlyTerm & Early Term. & Natural PG + LS & Episode termination & None \\
\bottomrule
\end{tabular}
\end{table}

\paragraph{Key design axes.} The algorithms differ along four orthogonal axes:
\begin{enumerate}
  \item \textbf{Policy update}: first-order (PG, PPO), or second-order natural gradient (NPG,
    TRPO) which provides better credit assignment at the cost of Fisher-vector-product computation.
  \item \textbf{Constraint enforcement}: Lagrangian relaxation, explicit quadratic programming
    (CPO), penalty functions, state augmentation, or primal feasibility switching (CRPO).
  \item \textbf{Multiplier adaptation}: fixed, gradient-based dual ascent, or PID control.
  \item \textbf{Environment modification}: standard CMDP or modified (Saute, Simmer,
    early termination).
\end{enumerate}

%%%%%%%%%%%%%%%%%%%%%%%%%%%%%%%%%%%%%%%%%%%%%%%%%%%%%%%%%%%%%
\section{Shared Implementation Details}
\label{sec:impl}
%%%%%%%%%%%%%%%%%%%%%%%%%%%%%%%%%%%%%%%%%%%%%%%%%%%%%%%%%%%%%

\subsection{Rollout and Data Collection}

All on-policy algorithms collect $T$ steps per epoch using \texttt{OnPolicyAdapter.rollout},
storing $(s_t, a_t, r_t, c_t, V^R(s_t), V^C(s_t), \log\piold(a_t|s_t))$ in the
\texttt{VectorOnPolicyBuffer}. At the end of each rollout, \texttt{finish\_path} is called
to compute GAE advantages (Eq.~\ref{eq:gae}) and value targets.

\subsection{Actor Network}

The policy is parameterised as a Gaussian distribution
$\pi_\theta(a|s) = \mathcal{N}(\mu_\theta(s),\,\mathrm{diag}(\sigma_\theta^2))$ with separate
mean and log-standard-deviation heads. For discrete action spaces, a categorical distribution
is used instead.

\subsection{Distributed Training}

OmniSafe uses MPI-style distributed training where gradients and statistics are averaged
across workers via \texttt{distributed.avg\_grads} and \texttt{distributed.dist\_avg}.

\subsection{Hyperparameter Defaults}

\begin{table}[h]
\centering
\caption{Common hyperparameter defaults in OmniSafe on-policy algorithms.}
\label{tab:hyperparams}
\small
\begin{tabular}{lll}
\toprule
\textbf{Parameter} & \textbf{Symbol} & \textbf{Default} \\
\midrule
Steps per epoch & $T$ & 20\,000 \\
Discount factor & $\gamma$ & 0.99 \\
GAE $\lambda$ (reward) & $\lambda_r$ & 0.95 \\
GAE $\lambda$ (cost) & $\lambda_c$ & 0.95 \\
Trust region budget & $\delta_{\mathrm{KL}}$ & 0.02 (PG/PPO), 0.01 (TRPO/CPO) \\
PPO clip radius & $\epsilon$ & 0.2 \\
Critic update iterations & --- & 40 \\
Mini-batch size & --- & 64 \\
Gradient norm clip & --- & 40.0 \\
CG iterations & --- & 15 \\
CG damping & $\alpha_{\mathrm{damp}}$ & 0.1 \\
Lagrange LR & $\eta_\lagm$ & varies \\
Cost limit & $d$ & 25 \\
\bottomrule
\end{tabular}
\end{table}

%%%%%%%%%%%%%%%%%%%%%%%%%%%%%%%%%%%%%%%%%%%%%%%%%%%%%%%%%%%%%
\bibliographystyle{plainnat}
\begin{thebibliography}{99}

\bibitem[Achiam et~al.(2017)]{achiam2017constrained}
Achiam, J., Held, D., Tamar, A., \& Abbeel, P. (2017).
\newblock Constrained policy optimization.
\newblock In \emph{International Conference on Machine Learning (ICML)}, pp. 22--31.

\bibitem[Altman(1999)]{altman1999constrained}
Altman, E. (1999).
\newblock \emph{Constrained Markov Decision Processes}.
\newblock CRC Press.

\bibitem[Chow et~al.(2017)]{chow2017risk}
Chow, Y., Ghavamzadeh, M., Janson, L., \& Pavone, M. (2017).
\newblock Risk-constrained reinforcement learning with percentile risk criteria.
\newblock \emph{Journal of Machine Learning Research}, 18(167), 1--51.

\bibitem[Ji et~al.(2023)]{ji2023omnisafe}
Ji, J., Zhang, B., Zhou, J., Pan, X., Dai, J., Li, P., Chen, J., Zheng, B., Liu, Y., \& Yang, Y. (2023).
\newblock OmniSafe: An infrastructure for accelerating safe reinforcement learning research.
\newblock \emph{arXiv preprint arXiv:2305.09304}.

\bibitem[Kakade(2001)]{kakade2001natural}
Kakade, S. (2001).
\newblock A natural policy gradient.
\newblock In \emph{Advances in Neural Information Processing Systems (NeurIPS)}, pp. 1531--1538.

\bibitem[Liu et~al.(2020)]{liu2020ipo}
Liu, Y., Ding, J., \& Liu, X. (2020).
\newblock IPO: Interior-point policy optimization under constraints.
\newblock In \emph{AAAI Conference on Artificial Intelligence}.

\bibitem[Liu et~al.(2022)]{liu2022penalized}
Liu, Z., Cen, Z., Isenbaev, V., Liu, W., Wu, S., Li, B., \& Zhao, D. (2022).
\newblock Constrained variational policy optimization for safe reinforcement learning.
\newblock In \emph{International Conference on Machine Learning (ICML)}.

\bibitem[Ray et~al.(2019)]{ray2019benchmarking}
Ray, A., Achiam, J., \& Amodei, D. (2019).
\newblock Benchmarking safe exploration in deep reinforcement learning.
\newblock \emph{arXiv preprint arXiv:1910.01708}.

\bibitem[Schulman et~al.(2015a)]{schulman2015high}
Schulman, J., Moritz, P., Levine, S., Jordan, M., \& Abbeel, P. (2015).
\newblock High-dimensional continuous control using generalized advantage estimation.
\newblock \emph{arXiv preprint arXiv:1506.02438}.

\bibitem[Schulman et~al.(2015b)]{schulman2015trust}
Schulman, J., Levine, S., Abbeel, P., Jordan, M., \& Moritz, P. (2015).
\newblock Trust region policy optimization.
\newblock In \emph{International Conference on Machine Learning (ICML)}, pp. 1889--1897.

\bibitem[Schulman et~al.(2017)]{schulman2017proximal}
Schulman, J., Wolski, F., Dhariwal, P., Radford, A., \& Klimov, O. (2017).
\newblock Proximal policy optimization algorithms.
\newblock \emph{arXiv preprint arXiv:1707.06347}.

\bibitem[Sootla et~al.(2022a)]{sootla2022saute}
Sootla, A., Cowen-Rivers, A.~I., Jafferjee, T., Wang, Z., Mguni, D.~H., Wang, J., \& Ammar, H. (2022).
\newblock Saut{\'e} RL: Almost surely safe reinforcement learning using state augmentation.
\newblock In \emph{International Conference on Machine Learning (ICML)}.

\bibitem[Sootla et~al.(2022b)]{sootla2022simmer}
Sootla, A., Cowen-Rivers, A.~I., Wang, J., \& Ammar, H.~B. (2022).
\newblock SIMMER: Adaptive budget scheduling for safe reinforcement learning.
\newblock In \emph{Advances in Neural Information Processing Systems (NeurIPS)}.

\bibitem[Stooke et~al.(2020)]{stooke2020responsive}
Stooke, A., Achiam, J., \& Abbeel, P. (2020).
\newblock Responsive safety in reinforcement learning by PID Lagrangian methods.
\newblock In \emph{International Conference on Machine Learning (ICML)}.

\bibitem[Sutton et~al.(1999)]{sutton1999policy}
Sutton, R.~S., McAllester, D., Singh, S., \& Mansour, Y. (1999).
\newblock Policy gradient methods for reinforcement learning with function approximation.
\newblock In \emph{Advances in Neural Information Processing Systems (NeurIPS)}, pp. 1057--1063.

\bibitem[Tessler et~al.(2018)]{tessler2018reward}
Tessler, C., Mankowitz, D.~J., \& Mannor, S. (2018).
\newblock Reward constrained policy optimization.
\newblock \emph{arXiv preprint arXiv:1805.11074}.

\bibitem[Xu et~al.(2021)]{xu2021crpo}
Xu, T., Liang, Y., \& Lan, G. (2021).
\newblock CRPO: A new approach for safe reinforcement learning with convergence guarantee.
\newblock In \emph{International Conference on Machine Learning (ICML)}.

\bibitem[Yang et~al.(2020)]{yang2020projection}
Yang, T.-Y., Rosca, J., Narasimhan, K., \& Ramadge, P.~J. (2020).
\newblock Projection-based constrained policy optimization.
\newblock In \emph{International Conference on Learning Representations (ICLR)}.

\bibitem[Yang et~al.(2022)]{yang2022constrained}
Yang, L., Ji, J., Dai, J., Zhang, L., Zhou, B., Li, P., Yang, Y., \& Pan, G. (2022).
\newblock Constrained update projection approach to safe policy optimization.
\newblock In \emph{Advances in Neural Information Processing Systems (NeurIPS)}.

\bibitem[Zhang et~al.(2020)]{zhang2020first}
Zhang, Y., Vuong, Q., \& Ross, K.~W. (2020).
\newblock First order constrained optimization in policy space.
\newblock In \emph{Advances in Neural Information Processing Systems (NeurIPS)}.

\end{thebibliography}

\end{document}
